% Generated draft: requires mathematical and source review.
\documentclass[11pt]{article}
\usepackage[margin=25mm]{geometry}
\usepackage{fontspec}
\setmainfont{Latin Modern Roman}
\usepackage{amsmath,amssymb,mathtools,unicode-math}
\setmathfont{Latin Modern Math}
\usepackage{enumitem,needspace,xurl,hyperref}
\setlength{\parindent}{0pt}
\setlength{\parskip}{5pt}
\setlength{\emergencystretch}{3em}
\newcommand{\N}{\mathbb{N}}
\newcommand{\Z}{\mathbb{Z}}
\newcommand{\Q}{\mathbb{Q}}
\newcommand{\R}{\mathbb{R}}
\newcommand{\C}{\mathbb{C}}
\newcommand{\F}{\mathbb{F}}
\newcommand{\A}{\mathbb{A}}
\newcommand{\PP}{\mathbb{P}}
\newcommand{\E}{\mathbb{E}}
\newcommand{\eps}{\varepsilon}
\newcommand{\dd}{\,\mathrm{d}}
\begin{document}
\section*{UnsolvedMath ID: 20002882}
\textbf{Source problem code:} catalogue entry
\subsection{Definitions and mathematical statement}
Let \(D=\{z\in\C^3:|z_i-c_i|<R_i,\ 1\le i\le3\}\) be an open polydisc with positive radii, and let
\[
\omega=\sum_{i=1}^3 a_i(z)\,\dd z_i
\]
be a holomorphic one-form on \(D\) satisfying the Frobenius identity \(\omega\wedge\dd\omega=0\). The approximation question attributed to J. E. Forn\ae ss asks whether there exists a sequence of polynomial one-forms
\[
h_\ell=\sum_{i=1}^3 p_{\ell,i}(z)\,\dd z_i,\qquad p_{\ell,i}\in\C[z_1,z_2,z_3],
\]
such that
\[
h_\ell\wedge\dd h_\ell=0\text{ identically on }\C^3,
\qquad
\max_i\sup_{z\in K}|p_{\ell,i}(z)-a_i(z)|\longrightarrow0
\]
for every compact \(K\subset D\). This is compact-open approximation of the coefficients; by holomorphy it also yields convergence of every derivative on compact subsets. Zeros of \(\omega\) and of the approximants are not excluded by the original problem.

A positive special case is reported in the supplied research aid. In centered coordinates \((x,y,z)\), suppose
\[
\omega=A(x,y,z)\,\dd x+B(x,y,z)\,\dd y,\qquad B\ne0\text{ everywhere on }D.
\]
Integrability forces \(r=A/B\) to be independent of \(z\). If \(B_n\) is the total-degree-\(n\) Taylor polynomial of \(B\), and \(r_n\) is that of \(r(x,y)\), then
\[
h_n=B_n(x,y,z)\bigl(r_n(x,y)\,\dd x+\dd y\bigr)
\]
is polynomial and exactly Frobenius-integrable, converges to \(\omega\) in compact-open \(C^\infty\), and satisfies \(j_0^nh_n=j_0^n\omega\). For each fixed compact \(K\subset D\), it is nonzero at every point of \(K\) for all sufficiently large \(n\). This eventual nonvanishing is a compactwise conclusion, not global nonvanishing of the polynomial form. The special case requires a global constant-tangent shape and a nowhere-zero coefficient on the entire polydisc; it is not the general local Frobenius theorem promoted to a global approximation result.

The aid also reports a broader rank-two pullback class and a reduction of a one-coefficient correction to a linear polynomial PDE, but provides no complete hypotheses or equations for these claims here. It warns that arbitrary Taylor truncation of all coefficients need not preserve Frobenius integrability. The general approximation question remains open according to the supplied aid, which mentions a broader restatement as Dinh--Sibony Problem 11.1. The source remark says a positive answer would lead to a nontrivial foliation of \(\C\PP^3\) by hypersurfaces; it does not specify the singularity or global construction conventions, and no nonsingular-foliation conclusion is asserted here.
\subsection{Short English statement}
Approximate an integrable holomorphic one-form on a polydisc by polynomial forms satisfying the Frobenius identity exactly. A reported special case admits factorized Taylor approximants with finite-jet matching and eventual compactwise nonvanishing, while the general problem remains open in the supplied report.
\subsection{Sources}
\begin{itemize}
\item J. E. Fornaess, Problem 1 and its remark, \emph{Approximation}, 2010 AIM workshop \emph{Emerging Applications of Complexity for CR Mappings}, page 3 according to the supplied verification, \url{https://aimath.org/WWN/crmappings/crmappings.pdf}.
\item The supplied research summary mentions Dinh--Sibony Problem 11.1 as a broader restatement, but gives no publication title, URL or full bibliographic entry.
\item Local UnsolvedMath registry, permanent ID 20002882, exact code \texttt{catalogue entry}; canonical location \texttt{aim-several-complex-variables-notes.json}, notes[39], dated 14 August 2026; statement labeled exact.
\item Research attempt 1: partial result, problem open at the run, rigor label proved, novelty confidence low; registry status partially solved. The special-case construction and current-status claim are supplied machine-generated research aids, not peer reviewed, requiring independent verification before citation.
\end{itemize}
\end{document}
